\begin{titlepage}
\centering
\vspace*{3cm}
{\LARGE Costas Arrays of Order 32\par}
\vspace{0.6cm}
{\Large Proceedings of an Investigation\par}
\vspace{2cm}
{\large Charles Gretton\textsuperscript{$\ast$}\par}
{\normalsize Australian National University\par}
\vspace{1cm}
{\normalsize October 2026\par}
\vfill
{\small\textsuperscript{$\ast$}Corresponding author:
\href{mailto:charles.gretton@anu.edu.au}{charles.gretton@anu.edu.au}\par}
\end{titlepage}

\chapter*{Licence and errata}
\addcontentsline{toc}{chapter}{Licence and errata}

\paragraph{Licence.} \textcopyright\ 2026 Charles Gretton. This work is licensed under the Creative Commons Attribution 4.0 International licence (CC BY 4.0).
To view a copy of the licence, visit \url{https://creativecommons.org/licenses/by/4.0/}. The software released
with it in the Costas-32 repository is licensed under the Apache License, Version 2.0
(\url{https://www.apache.org/licenses/LICENSE-2.0}). Material by others that is cited or reproduced here,
including published data and catalogues, remains under its original terms.

\paragraph{Errata.} We welcome reports of errors, omissions and overlooked prior work. Please send them to the
corresponding author, Charles Gretton, at \href{mailto:charles.gretton@anu.edu.au}{charles.gretton@anu.edu.au}.
Every report that leads to a correction will be acknowledged in the errata and in any revised edition,
cited as a private communication; where the point was already made in a publication that appeared before
this report's publication date, that publication will be cited as well.

\chapter*{AI usage acknowledgment}
\addcontentsline{toc}{chapter}{AI usage acknowledgment}

Large language models were used extensively throughout this investigation and in preparing this
volume. Claude Opus 5.5 and GPT 5.6 SOL were used for the report and for the proceedings as a
whole. Their use covered developing and checking mathematical arguments; designing, writing,
testing and running the software; planning, monitoring and analysing the computational campaigns;
searching and assessing the literature; and drafting, reviewing and revising the text, including
assembling these proceedings from the project's reports.

The author directed the work, decided what to run and what to claim, and checked its major
claims. He has not made a complete inventory of the claims in this volume, nor checked each one
individually; doing so would take far longer than the investigation itself, so errors may remain,
and reports of them are welcome (see the errata above). What the volume provides instead is the
means to check: each computational claim rests on released software and recorded outputs, and the
report states how its claims were checked. The use of models also shaped the methods: much of the
emphasis on fail-closed accounting, independent verification of every answer and hypotheses stated
before the runs, described in the note on software engineering, exists so that results produced with heavy machine assistance can be audited without
trusting the process that produced them.

The author calls a release of this kind a \emph{virtual proceedings}: a self-contained manuscript
published together with the code package that produced its results --- the volume and its source,
the software needed to rerun every experiment, and the derived evidence behind every reported
number --- as one open repository rather than through a conventional venue. The name reflects both
how such work is now produced, largely in dialogue with AI tools under the author's direction, and
how it is meant to be used: read, rerun and audited by others.

\chapter*{Preface}
\addcontentsline{toc}{chapter}{Preface}

Orders 32 and 33 are the smallest orders for which the existence of a Costas array is unknown. This
volume collects the results of a systematic investigation of order 32 as ten self-contained
papers, one for each strand of the work. Each paper is readable on its own at the level of its main claim. Shared notation, recurring definitions and release conventions are collected once, in the front matter, and each paper states its dependencies on the others.

The investigation does not resolve existence at order 32. What it does is close a series of
structured routes by which an order-32 array could arise, and measure or explain why others fail.
The table summarises the ten papers and the standing of their results.

\begin{center}
\small
\begin{tabular}{@{}p{0.04\textwidth}>{\raggedright\arraybackslash}p{0.42\textwidth}>{\raggedright\arraybackslash}p{0.44\textwidth}@{}}
\toprule
& result & standing \\
\midrule
1 & No known construction reaches order 32; no circular Costas array of order 32 exists. &
Mostly from the literature; single periodicity remains open pending the Golomb--Taylor conjecture. \\
2 & No order-32 array is a two-point augmentation of the 664 generated algebraic order-30 cores;
the closest miss has a complete genealogy. & Computational; relative to the generated order-30
catalogue. \\
3 & No order-32 array contains an order-29 or an order-28 Costas deletion minor. &
Computational; quantifies over complete published censuses; neither result is known to imply the other. \\
4 & The order-29 census, reconstructed and independently verified. & Verification, partly
kernel-checked in Lean; completeness rests on the published enumeration. \\
5 & A telescoping-sum propagation rule, proved sound. & Passes a frozen benchmark; slows exact search
in practice. \\
6 & A first-moment model forecasting about $10^{-4}$ sporadic classes at order 32. & Conditional
forecast; an equally good fit diverges. \\
7 & At least 112 of the 240 diagonal second differences of any order-32 permutation are even. &
Proved; does not bear on existence. \\
8 & No transpose-symmetric Costas array of order 32 exists. & Known since 1993; determined
independently here twice, by SAT and by exhaustive direct search. \\
9 & An exhaustive search of all order-32 permutations, by a straightforward method, would take about $1.5\times10^6$ core-years as first written, about $1.9\times10^5$ once engineered. & Measured, with validated estimators; the faster 2011 method would need an estimated $10^4$ core-years. \\
10 & Late dead prefixes are refuted by a pair-marginal linear relaxation; nothing tested refutes earlier prefixes more cheaply than search. & Sampled and independently replicated; caps any rule acting at the transition at about nine-fold at order 32. \\
\bottomrule
\end{tabular}
\end{center}

Two cautions apply throughout. First, most results here are computational: they rest on the
correctness of the software, which is released, tested and in places independently re-implemented,
but not formally certified except where a paper says so. Each paper ends with a section stating
the standing of its claims. Second, the result of Paper~8 turned out to have been published in 1993
and confirmed in 2009, which the project discovered only after its own campaign was complete. The
other strands have not yet been checked against the literature as thoroughly, and their novelty
should be read with that in mind.

Some papers point to the project's consolidated report for material outside their scope, such as
structural approaches that were tried and refuted.

\paragraph{Data and code availability.} The consolidated report, all software, Lean sources,
catalogues, logs and result ledgers referred to in this volume are released together as one software
archive (release \texttt{costas32-1.3}), whose manifest records the hash of every file. A persistent
identifier for the archive is to be assigned on deposit.

\paragraph{Dependencies and standing.} The table records, for each paper, the kind of result, what it
depends on, and what it takes from outside the project.

\begin{center}
\small
\begin{tabular}{@{}p{0.04\textwidth}>{\raggedright\arraybackslash}p{0.21\textwidth}>{\raggedright\arraybackslash}p{0.23\textwidth}>{\raggedright\arraybackslash}p{0.23\textwidth}>{\raggedright\arraybackslash}p{0.19\textwidth}@{}}
\toprule
& result type & depends on & external premise & standing \\
\midrule
1 & background and analysis & shared notation & published theorems on constructions and periodicity & cited or derived \\
2 & catalogue-relative exclusion & catalogue generator & none; order 30 is unenumerated & computational, uncertified \\
3 & census-complete exclusions & rung engine; Paper~4 & completeness of the order-28 and order-29 censuses & computational, uncertified \\
4 & census verification & published representatives & completeness of the order-29 census & partly kernel-checked \\
5 & inference rule and evaluation & none & none & proved rule; measured performance \\
6 & conditional model & Paper~1; published counts & the model form & heuristic \\
7 & theorem and bounds & none & none & proved; computational \\
8 & independent confirmation & regional encoding, solvers, audit & none; the 1993 and 2009 counts are for comparison & computational, uncertified \\
9 & cost measurement & Papers~1 and~8 & the reported cost of the 2011 enumeration & estimated, validated \\
10 & obstruction search & Papers~2 and~9 & the complete catalogues of orders 12--15 & sampled, replicated \\
\bottomrule
\end{tabular}
\end{center}

What remains open is the residual problem, stated more fully in the closing chapter: order-32 arrays that are not algebraic, not periodic,
not geometrically symmetric, contain no Costas deletion minor of order 28 or 29, and are not
augmentations of algebraic order-30 cores. No structural classification in this volume controls
that class, and exhaustive enumeration of all order-32 permutations remains far out of reach.


\chapter*{A note on software engineering}
\addcontentsline{toc}{chapter}{A note on software engineering}

The investigation was run as a software project as much as a mathematical one, and several terms
from that practice recur in the papers. Two of them, \emph{rungs} and \emph{gates}, carry technical
weight and are introduced here.

\paragraph{Rungs.} Several searches are organised as a ladder of finite exhaustive searches indexed
by a parameter. Each \emph{rung} is one complete run, with its own declared coverage and its own
verdict. The main instance is the augmentation ladder of Papers~2 and~3. Rung $k$ searches every
permutation of order $n$ obtained from a Costas core of order $n-k$ by inserting $k$ columns, $k$
rows and $k$ dots. Rung~1 runs from order 30 to order 31; rungs 2, 3 and 4 target order 32 from
cores of orders 30, 29 and 28. Rungs are run in
increasing cost, and each is self-contained: a negative result at one rung is a theorem about that
rung alone. Rungs need not imply one another; Paper~3 explains why neither of rungs 3 and 4 is known to imply the other.

\paragraph{Gates.} A \emph{gate} is an acceptance criterion written down and frozen before the run
it judges, in the manner of a pre-registered experiment or a release gate in continuous
integration. Freezing it prevents the criterion from being tuned to the result. Paper~5 evaluates a
propagation rule in two stages: Stage~1 on a benchmark of partial assignments, and Stage~2 inside an
exact search. It uses two kinds of gate. A \emph{scientific gate} asks whether the rule deduces
more than the baseline without ever rejecting a completable assignment. An \emph{operational gate}
asks whether it makes search faster. A result can pass one and fail the other, and Paper~5's does. (Paper~8 uses ``Stage~1'' independently, for the first stage of its own campaign.)

\paragraph{Related practice.} Other recurring terms are as follows.
\begin{itemize}
\item \emph{Fail-closed accounting.} A run counts as complete only when its declared coverage is known to cover the whole search space, by a proof such as Paper~8's partition and cover theorems or by construction, as with Paper~3's one work unit per dihedral orbit, and every unit of that coverage reports an explicit terminal status. Anything missing, failed or unparsed counts as not done.
\item \emph{Census match.} Before a census-complete search runs, its catalogue is checked against
the published census figures and its hash recorded.
\item \emph{Work units and independent merging.} Large searches are split into deterministic units
whose outputs are merged and checked by a separate tool.
\item \emph{Independent re-scoring.} Every reported near-solution is re-scored by a separately
written scorer.
\item \emph{Release gate.} A repository self-test that must pass before any archive is released.
\end{itemize}

\chapter*{Notation and glossary}
\addcontentsline{toc}{chapter}{Notation and glossary}

Papers 1--7 write a permutation as a sequence of heights $h$, with the dot of column $i$ in row
$h_i$; Paper~8 writes it as $\pi$. All indices are zero-based.

\section*{Notation}
\begin{longtable}{@{}l >{\raggedright\arraybackslash}p{0.62\textwidth}@{}}
\toprule
Symbol & Meaning \\
\midrule
\endhead

$n$ & Order of a Costas array; the board is $n\times n$. Usually $n=32$. \\
$h$ & A permutation $h:\{0,\dots,n-1\}\to\{0,\dots,n-1\}$, written zero-based.
      $h_i$ is the \emph{height} (row) of the dot in column $i$. \\
$\Delta_{m,i}$ & $h_{i+m}-h_i$, the vertical difference at horizontal
      separation $m$ starting at column $i$. \\
$m$, $d$ & Horizontal separation (distance) between two columns. $d$ is used in
      solver contexts, $m$ in partition-counting contexts; they denote the same
      quantity. \\
$e$ & A signed vertical difference; $e\in\mathcal V$. \\
$g$ & An absolute vertical distance, $g=|e|$. \\
$\mathcal V$ & The signed palette $\{-(n-1),\dots,-1,1,\dots,n-1\}$,
      $|\mathcal V| = 2n-2$; for $n=32$, $|\mathcal V|=62$. \\
$L_m$ & $n-m$, the number of column pairs at separation $m$. \\
$\score(h)$ & Costas defect: the number of repeated $(d,\Delta)$ pairs,
      $\sum_{d,e}\max(0,\chi_{d,e}-1)$, over $1\le d\le n-1$ and nonzero $e$ with $|e|\le n-1$ (differences of a permutation are never zero). $\score(h)=0$ iff $h$ is Costas. \\
$\chi_{m,e}$ & $\#\{i: \Delta_{m,i}=e\}$, the occupancy of signed value $e$
      at separation $m$. \\
$y_{m,g}$ & $\chi_{m,g}+\chi_{m,-g}\in\{0,1,2\}$, magnitude occupancy. \\
$T_m$ & $\sum_i \Delta_{m,i}$, the telescoping row sum at separation $m$. \\
$N_q(r)$ & $\#\{e\in\mathcal V : e\equiv r \pmod q\}$, the signed residue
      capacity. \\
$\vtwo(x)$ & The 2-adic valuation of $x$; $\vtwo(0)=\infty$. \\
$D_i$ & Domain of the height variable $h_i$ during propagation. \\
$E_{i,j}$, $D_{ij}$ & Domain of the edge (difference) variable $c_{ij}=h_j-h_i$. \\
$k$ & Number of columns and rows deleted in an augmentation rung; the core has
      order $n-k$. \\
$K$ & A Costas core array of order $n-k$. \\
$C$, $R$ & The deleted column set and deleted row set, each of size $k$. \\
$\delta,\epsilon$ & Compressed displacement of a core pair: column distance and
      row difference \emph{within the core}. \\
$D_4$ & The dihedral group of order 8 acting on the square; the Costas
      symmetry group. \\
$W_1(p,g,c)$ & Welch construction, order $p-1$: $h_i = g^{i+c}\bmod p - 1$. \\
$W_2,W_3$ & $W_1$ with one or two corner dots deleted, orders $p-2$, $p-3$. \\
$G_2(q,\alpha,\beta)$ & Golomb construction over $\F_q$, order $q-2$:
      dot at $(i,j)$ iff $\alpha^i+\beta^j=1$. \\
$G_3,G_4,G_5$ & $G_2$ with corner dots deleted, orders $q-3,q-4,q-5$. \\
$C_{31}$ & The unique order-31 Costas array reachable by one-dot augmentation
      of the algebraic order-30 cores (one $D_4$ orbit). \\
$N_{32}$ & The known order-32 permutation of defect 2; the corner extension of
      $C_{31}$. \\
$\pi$ & In Paper~8, a permutation, written $\pi$ rather than $h$; transpose symmetry is
      $\pi=\pi^{-1}$. \\
$\Fix(\pi)$, $f$ & The fixed points of an involution $\pi$, and their number. \\
$F$, $R_F$ & A fixed-point set, and the region of Costas involutions whose fixed-point set is
      exactly $F$ (Paper~8). \\
$x_{ab}$, $h_{ij}$, $p_{s,x,\delta}$ & The edge, image and displacement-indicator variables of
      Paper~8's encoding of a region. \\
$Z(\pi)$ & In Paper~8, the number of large-large pairs: pairs $\{x,y\}$ with $|x-y|\ge16$ and
      $|\pi(x)-\pi(y)|\ge16$. \\
$k$ (Paper~8) & The number of split coordinates in a native split. In Papers~2 and~3, $k$ is the
      rung parameter instead. \\
$M(D,k)$ & In Paper~8, the number of messages of a native split with $k$ split coordinates and $D$
      free partners. \\
\bottomrule
\end{longtable}

\section*{Glossary}
\begin{description}[style=nextline,leftmargin=0pt,itemsep=2pt]
\item[Costas array] A permutation array of order $n$ in which all $\binom{n}{2}$ vectors connecting
pairs of dots are distinct. Equivalently, for every horizontal separation $m$, the vertical
differences $h_{i+m}-h_i$ are pairwise distinct.
\item[Displacement] The vector $(j-i,\,h_j-h_i)$ between the dots of columns $i<j$.
\item[Defect, score] The number of repeated displacement vectors, $\score(h)$. A Costas array has
defect 0; the known order-32 near-solution $N_{32}$ has defect 2.
\item[Dihedral orbit, equivalence class] The set of arrays obtained from one array under the eight
symmetries of the square, $D_4$. Censuses count arrays both individually and by class.
\item[Census] A complete list of the Costas arrays of one order, obtained by exhaustive enumeration.
Complete censuses exist through order 29.
\item[Census-complete, catalogue-relative] A result is census-complete if it quantifies over a
complete census, and catalogue-relative if it quantifies only over a generated catalogue, such as
the 664 order-30 arrays of Paper~2, at an order that has not been enumerated.
\item[Welch and Golomb constructions] The two principal algebraic constructions of Costas arrays,
over prime fields and finite fields respectively; see the notation table.
\item[Circular and singly periodic sequences] A Costas sequence $a_0,\dots,a_{n-1}$ with values in
$\{1,\dots,n\}$ is \emph{circular} if, for every lag $k$, the circular differences
$a_{(i+k)\bmod n}-a_i$ are distinct modulo $n+1$: its dots, read in $\Z_n\times\Z_{n+1}$, have
distinct displacements. It is \emph{singly periodic} if every cyclic shift of its terms is again
Costas. Circularity lies between single periodicity and Moreno's shifting property. Circular
sequences exist only when $n+1$ is prime (Golomb and Moreno, 1996) and are exactly the Welch
sequences (Muratovi\'c-Ribi\'c, Pott, Thomson and Wang, 2015); every circular sequence is an
ordinary Costas array, but not conversely. Arrays required to stay Costas when extended periodically
in both directions are a further notion, studied separately. See Paper~1.
\item[Corner deletion, corner addition] Removing or adding a dot at a corner of the board, together
with its row and column. The compression is then a translation, so these are the deletions that
preserve the Costas property in general.
\item[Sporadic array] A Costas array not accounted for by a recognised generated or emergent
construction. The category shrinks whenever a new construction is found. Order 28 is the first
order above 5 with no sporadic arrays, and order 29 has none either.
\item[Emergent array] An array produced by a semi-empirical construction rather than an algebraic
family, such as the 16 order-29 arrays produced by the Rickard--Welch operation.
\item[Augmentation rung] The exhaustive search over all permutations of order $n$ formed by placing
a Costas core of order $n-k$ in a compressed sub-board and adding $k$ dots; see the note on software
engineering.
\item[Deletion minor] The permutation array left after deleting $k$ columns, $k$ rows and the $k$
dots they share from a larger permutation array, and closing up the gaps: the inverse of a rung-$k$
augmentation.
\item[Core-first pruning] In a rung search, testing the placed core's own displacements before
considering the added dots. Exact for defect 0, since a repeated displacement among a subset of the
dots is a repeated displacement of the whole.
\item[Involution, fixed point, transposition] A permutation equal to its inverse; a coordinate $c$
with $\pi(c)=c$, whose dot lies on the main diagonal; a pair $(a\,b)$ with $\pi(a)=b$ and $\pi(b)=a$.
\item[Geometric and transpose symmetry] An array is geometrically symmetric if a non-identity motion
of the square fixes it, and transpose-symmetric if the main-diagonal reflection does. For orders of
at least 3 the two notions agree up to the choice of diagonal (Paper~8). The literature calls
transpose-symmetric arrays \emph{diagonal}.
\item[Anti-reflective and consecutive arrays] Symmetry types studied in the literature that are not
motions of the square. They are mentioned for comparison only.
\item[Levels of assurance] A \emph{validated} computation agrees with independent implementations,
regression tests or known answers; an \emph{audited} computation has its coverage and execution
records independently checked; a \emph{machine-checked witness} is an object, such as a Costas
permutation or a Lean theorem, checked independently; a \emph{proof-certified} unsatisfiable verdict
has a checked DRAT or LRAT proof; an \emph{uncertified} one rests on the solver's exhaustion.
\item[Region] In Paper~8, the set of Costas involutions with a given fixed-point set.
\item[Native split, message, unsplit branch] Paper~8's distributed decomposition fixes the partners
of $k$ chosen coordinates; each assignment of partners is a message, and a message in which two
chosen coordinates are paired with each other is an unsplit branch.
\item[Seed-refuted] In Paper~8, a pairing that repeats a displacement with the fixed dots alone.
\item[Case] In Paper~8, a region with one transposition fixed.
\item[Gate] An acceptance criterion frozen before the run it judges; see the note on software
engineering.
\item[CDCL, DRAT and LRAT] Conflict-driven clause learning, the algorithm of modern SAT solvers; and
machine-checkable proof formats for unsatisfiability. Without such a proof, an unsatisfiable verdict
is a computational observation rather than a certificate.
\item[Verdicts] Paper~8 records \st{SAT-VALIDATED} (a model that passed independent checking),
\st{UNSAT-UNCERTIFIED} (normal solver exhaustion, no checked proof), \st{TIMEOUT} and \st{ERROR}. A
region closed by an \st{UNSAT-UNCERTIFIED} verdict is \emph{search-closed}.
\end{description}


\tableofcontents
